\lecture{8.7}{Sequential Regression G-Computation}
\label{lec:8-7}

\subsection{Backward Recursive Regressions}

Sequential regression is an implementation of the longitudinal g-computation
formula that avoids explicitly estimating the full conditional distribution of
each time-varying covariate. Instead, it works backward from the final outcome.

For $K=1$ and regime $(a_0,a_1)$:
\[
  \Qbar_2(L_0,L_1)=\E(Y\given A_0=a_0,A_1=a_1,L_0,L_1),
\]
\[
  \Qbar_1(L_0)=\E\{\Qbar_2(L_0,L_1)\given A_0=a_0,L_0\},
\]
\[
  \psi(a_0,a_1)=\E\{\Qbar_1(L_0)\}.
\]
Sequential regression estimates $\Qbar_2$, predicts it for all individuals, then
uses those predictions as the outcome in a regression for $\Qbar_1$ among
individuals with $A_0=a_0$.

\begin{figure}[H]
\centering
\begin{tikzpicture}[
  node distance=.8cm,
  box/.style={draw, rounded corners=2pt, minimum width=3.5cm, minimum height=.74cm},
  arr/.style={-{Latex[length=2.2mm]}, thick}
]
\node[box] (s1) {Fit final outcome regression\\ among $\Abar_K=\abar_K$};
\node[box, below=of s1] (s2) {Predict pseudo-outcome\\ for all histories};
\node[box, below=of s2] (s3) {Regress pseudo-outcome\\ among $\Abar_{K-1}=\abar_{K-1}$};
\node[box, below=of s3] (s4) {Repeat backward and\\ average $\Qbar_1(L_0)$};
\draw[arr] (s1) -- (s2);
\draw[arr] (s2) -- (s3);
\draw[arr] (s3) -- (s4);
\end{tikzpicture}
\caption{Sequential regression implements g-computation by backward recursion.}
\label{fig:ch8-sequential-regression}
\end{figure}

\subsection{Pseudo-Outcomes}

A pseudo-outcome is a fitted value from a later regression that becomes the
outcome for an earlier regression. In the two-timepoint case:
\begin{enumerate}
  \item Fit a regression for $Y$ among those with $A_0=a_0,A_1=a_1$.
  \item Predict $\Qbar_{2,n}(L_{0i},L_{1i})$ for every individual.
  \item Among those with $A_0=a_0$, regress $\Qbar_{2,n}(L_{0i},L_{1i})$ on
  $L_{0i}$.
  \item Predict $\Qbar_{1,n}(L_{0i})$ for every individual.
  \item Average these predictions.
\end{enumerate}

This is a regression-based way to perform the integration over $L_1$ required by
the g-computation formula.

\subsection{Binary Time-Varying Confounders}

If $L_1$ is binary or low-dimensional discrete, one can also implement
g-computation by explicitly summing over possible $L_1$ values:
\[
  \Qbar_1(L_0)
  =
  \sum_{l_1}
  \Qbar_2(L_0,l_1)
  \Prob(L_1=l_1\given A_0=a_0,L_0).
\]
For binary $L_1$,
\[
\begin{aligned}
  \Qbar_1(L_0)
  &=
  \Qbar_2(L_0,1)\Prob(L_1=1\given A_0=a_0,L_0)\\
  &\quad+
  \Qbar_2(L_0,0)\Prob(L_1=0\given A_0=a_0,L_0).
\end{aligned}
\]
This approach is transparent, but it becomes difficult when time-varying
covariates are continuous or high-dimensional.

\subsection{Estimating $\Qbar_{K+1},\Qbar_K,\ldots,\Qbar_1$}

For a general fixed regime $\abar$, sequential regression proceeds as follows:
\begin{estimationbox}
\begin{enumerate}
  \item Estimate
  \[
    \Qbar_{K+1}(\Lbar_K)=
    \E(Y\given \Abar_K=\abar_K,\Lbar_K)
  \]
  using observations with $\Abar_K=\abar_K$.
  \item Predict $\Qbar_{K+1,n}(\Lbar_{K,i})$ for all individuals.
  \item For $k=K,K-1,\ldots,1$, estimate
  \[
    \Qbar_k(\Lbar_{k-1})=
    \E\{\Qbar_{k+1}(\Lbar_k)\given
    \Abar_{k-1}=\abar_{k-1},\Lbar_{k-1}\}
  \]
  by regressing the current pseudo-outcome on $\Lbar_{k-1}$ among those with
  $\Abar_{k-1}=\abar_{k-1}$.
  \item Estimate
  \[
    \psi(\abar)=\frac{1}{n}\sum_{i=1}^n \Qbar_{1,n}(L_{0i}).
  \]
\end{enumerate}
\end{estimationbox}

\subsection{R Implementation}

For two treatment times and regime $(1,1)$:

\begin{lstlisting}[language=R]
full_data <- data.frame(L0 = L0, A0 = A0, L1 = L1, A1 = A1, Y = Y)

# Step 1: final outcome regression among regime followers
data_11 <- subset(full_data, A0 == 1 & A1 == 1)
fit_Q2 <- glm(Y ~ L0 + L1, data = data_11)

# Step 2: predict Q2 for everyone at their observed L0, L1
full_data$Q2n <- predict(fit_Q2, newdata = full_data, type = "response")

# Step 3: regress the pseudo-outcome among those with A0 = 1
data_1 <- subset(full_data, A0 == 1)
fit_Q1 <- glm(Q2n ~ L0, data = data_1)

# Step 4: predict Q1 for everyone and average
full_data$Q1n <- predict(fit_Q1, newdata = full_data, type = "response")
psi_11 <- mean(full_data$Q1n)
\end{lstlisting}

If the outcome is binary, logistic regression may be used for both steps:

\begin{lstlisting}[language=R]
fit_Q2 <- glm(Y ~ L0 + L1, data = data_11, family = binomial())
full_data$Q2n <- predict(fit_Q2, newdata = full_data, type = "response")

fit_Q1 <- glm(Q2n ~ L0, data = data_1, family = binomial())
full_data$Q1n <- predict(fit_Q1, newdata = full_data, type = "response")
psi_11 <- mean(full_data$Q1n)
\end{lstlisting}

\subsection{Comparing With IPTW}

Sequential regression and IPTW estimate the same counterfactual mean under the
same identification assumptions, but they rely on different nuisance functions.
Sequential regression relies on the outcome regression sequence
$\Qbar_{K+1},\ldots,\Qbar_1$. IPTW relies on the treatment mechanism sequence
$g_0,\ldots,g_K$.

Sequential regression often has lower variance because it does not use inverse
probability weights. However, it can be biased if the sequential outcome
regressions are misspecified. IPTW may be more robust to outcome model
misspecification but can be unstable under near-positivity violations.

\begin{keyidea}
Sequential regression turns nested expectations into a sequence of prediction
problems. At each step, the outcome is the predicted counterfactual mean from
the next later step.
\end{keyidea}
